\chapter{Options multi-actifs}\label{ch:dv:multi-asset-options}

Une note verse un coupon en fin d'année seulement si ses trois actions valent alors au moins 70~\% de leur
niveau de départ. Le client voit trois noms familiers et une condition qui semble facile à remplir. Le dealer voit une
position en corrélation qu'aucun marché coté n'affiche. Si les trois actions étaient indépendantes, avec des volatilités de
25~\%, 30~\% et 35~\%, la probabilité de verser le coupon serait de 62~\%. Si elles évoluaient ensemble, elle serait de 79~\%. Le
dealer qui a vendu le coupon est short de cet écart. Ce chapitre évalue des options sur plusieurs actifs, sur paniers,
\omterm{def:dv:multi-asset-options:worstof}{options worst-of} et best-of, et montre comment les desks lisent la corrélation à partir des marchés d'options sur indice et sur actions individuelles. Il
traite ensuite des \omterm{def:dv:multi-asset-options:dispersion}{trades de dispersion}, portant sur l'écart entre les deux, de la \omterm{def:dv:multi-asset-options:local}{corrélation locale}, qui s'ajuste au skew d'indice, et les
\omterm{def:dv:multi-asset-options:quanto}{options quanto} et composites qui font entrer une seconde devise dans le payoff.

\section{Corrélation et paniers}

Tous les modèles de ce chapitre partent d'actifs lognormaux corrélés~:
$dS_i/S_i=\mu_i\,dt+\sigma_i\,dW^i$ avec $d\langle W^i,W^j\rangle=\rho_{ij}\,dt$. Les trajectoires sont simulées par la
factorisation de Cholesky $L$ de la matrice de corrélation, $LL^\top=\rho$ (livre~4, chapitre~25), appliquée à des
normales indépendantes, puis on évalue par Monte-Carlo (livre~4, chapitre~26). La construction retient chaque $\sigma_i$ comme un nombre ou comme une
fonction de \omterm{def:dv:local-volatility:model}{volatilité locale} $\sigma_i(t,S_i)$, de sorte que le smile propre à chaque actif soit respecté.

\begin{definition}[Option sur panier]\label{def:dv:multi-asset-options:basket}
Une \emph{option sur panier}\index{option sur panier} est une option sur une somme pondérée d'actifs, typiquement de leurs
performances~: $\bigl(\sum_iw_iS_{i,T}/S_{i,0}-K\bigr)^+$ pour un call.
\end{definition}

La variance d'un panier vaut $\sum_{ij}w_iw_j\rho_{ij}\sigma_i\sigma_j$ dans l'approximation lognormale. Elle est inférieure à la
moyenne pondérée des variances dès que les corrélations sont inférieures à un, ce qui est la diversification que le
client achète. Un panier de lognormales n'est pas lognormal, mais une lognormale aux deux mêmes premiers moments l'évalue bien. On égale
$M_1=\sum w_iF_i$ et $M_2=\sum_{ij}w_iw_jF_iF_je^{\rho_{ij}\sigma_i\sigma_jT}$, puis on évalue avec Black à la volatilité
$\sqrt{\ln(M_2/M_1^2)/T}$. Sur un panier à pondérations égales des trois actions, à un an, l'approximation est à moins de
0.0004 du prix à la monnaie simulé pour toute corrélation comprise entre 0 et 0.95.

\begin{omfigure}
\begin{tikzpicture}
\begin{groupplot}[group style={group size=2 by 1,horizontal sep=1.6cm},
  width=6.6cm,height=5cm,
  tick label style={font=\small},label style={font=\small},grid=major,grid style={gray!25},table/col sep=comma]
\nextgroupplot[xlabel={corrélation},ylabel={prix (par unité de notionnel)},xmin=0,xmax=1,ymin=0,ymax=0.3,
  yticklabel style={/pgf/number format/fixed,/pgf/number format/precision=2},
  title={\small options à la monnaie, à un an},
  legend columns=1,legend cell align=left,legend style={font=\footnotesize,at={(0.5,-0.32)},anchor=north,draw=none,fill=none}]
\addplot[omThm,very thick,mark=*,mark size=1.4pt] table[x=rho,y=worst_put]{figdata/derivatives/17-multi-asset-options/three_share.csv};
\addplot[omDef,very thick,mark=square*,mark size=1.4pt] table[x=rho,y=best_call]{figdata/derivatives/17-multi-asset-options/three_share.csv};
\addplot[omProp,very thick,mark=triangle*,mark size=1.6pt] table[x=rho,y=basket_call]{figdata/derivatives/17-multi-asset-options/three_share.csv};
\legend{put worst-of,call best-of,call sur panier}
\nextgroupplot[xlabel={corrélation},ylabel={probabilité},xmin=0,xmax=1,ymin=0.55,ymax=0.85,
  title={\small les trois au-dessus de 70~\%}]
\addplot[omThm,very thick,mark=*,mark size=1.4pt] table[x=rho,y=worst_digital]{figdata/derivatives/17-multi-asset-options/three_share.csv};
\end{groupplot}
\end{tikzpicture}

\omcaption{Trois actions de volatilités 25~\%, 30~\% et 35~\%, à un an, taux nuls, corrélation deux à deux égale.
À gauche~: les puts worst-of et les calls best-of perdent de la valeur à mesure que la corrélation monte, et les calls sur panier en gagnent. À droite~: la
probabilité que les trois actions terminent au-dessus de 70~\% de leur niveau de départ, soit la condition de coupon de la note. Données~: celles du
tutoriel.}\label{fig:dv:multi-asset-options:correlation}
\end{omfigure}

\section{Worst-of et best-of}

\begin{definition}[Option worst-of]\label{def:dv:multi-asset-options:worstof}
Une \emph{option worst-of}\index{option worst-of} paie sur la moins bonne performance parmi plusieurs actifs, par exemple
$(K-\min_iS_{i,T}/S_{i,0})^+$ pour un put worst-of.
\end{definition}

\begin{definition}[Option best-of]\label{def:dv:multi-asset-options:bestof}
Une \emph{option best-of}\index{option best-of} paie sur la meilleure performance, par exemple
$(\max_iS_{i,T}/S_{i,0}-K)^+$ pour un call best-of.
\end{definition}

La corrélation agit sur ces payoffs dans le sens inverse de celui d'un panier. Lorsque les actifs évoluent indépendamment, la
moins bonne performance se situe bien en dessous de la moyenne, et la meilleure bien au-dessus. Lorsqu'ils évoluent ensemble, toutes les performances sont
proches, et les \omterm{def:dv:multi-asset-options:worstof}{options worst-of} et best-of ressemblent à des options sur un seul actif. Ainsi, un put worst-of et un call best-of perdent
de la valeur à mesure que la corrélation monte, et un call sur panier en gagne
(\cref{fig:dv:multi-asset-options:correlation}).

\begin{example}[Les trois actions de la note]\label{ex:dv:multi-asset-options:three}
Pour les trois actions (volatilités 25~\%, 30~\% et 35~\%, à un an, taux nuls), le put worst-of à la monnaie vaut
0.250 par unité de notionnel avec des actions indépendantes et 0.152 pour une corrélation de 0.95. Le call best-of baisse de 0.284 à
0.160, et le call sur panier monte de 0.071 à 0.118. La probabilité que les trois actions terminent à 70~\% au moins de leur niveau de départ
passe de 0.616 à 0.787. Pour une corrélation de 0.4, elle vaut 0.673, et 0.707 pour 0.6~: une erreur de cinq points
de corrélation déplace la valeur du coupon d'environ un point de probabilité.
\end{example}

Une note qui verse un coupon lorsque toutes les actions sont au-dessus d'un niveau, et rembourse moins que le capital lorsque la moins bonne
action est tombée sous un niveau plus bas (les autocallables du chapitre~18), laisse le dealer long d'un put worst-of et short
d'un digital worst-of. Les deux perdent de la valeur quand la corrélation monte, donc le dealer est short de corrélation, et le marché des options sur indice est
où cette exposition peut être en partie compensée. Les autres expositions du dealer sont les cross-gammas, soit les dérivées secondes par
paires d'actifs (livre~6, chapitre~3). Elles font dépendre la couverture de chaque action des mouvements des autres, et
elles sont les plus fortes près de la barrière où le titre le moins performant change.

\section{Corrélation implicite, dispersion et skew de corrélation}

Un indice est un panier de ses composantes, si bien que les options sur indice et les options sur actions individuelles impliquent ensemble une corrélation.

\begin{definition}[Corrélation implicite]\label{def:dv:multi-asset-options:implied}
La \emph{corrélation implicite}\index{corrélation implicite} d'un indice aux pondérations $w_i$ est l'unique corrélation
$\bar\rho$ qui rend cohérentes les volatilités implicites $\sigma_i$ des composantes avec la volatilité implicite de l'indice
$\sigma_I$~:
\[
\bar\rho=\frac{\sigma_I^2-\sum_iw_i^2\sigma_i^2}{\bigl(\sum_iw_i\sigma_i\bigr)^2-\sum_iw_i^2\sigma_i^2}.
\]
\end{definition}

Calculée avec des volatilités à la monnaie, elle est un nombre unique. Calculée à chaque strike, en utilisant la volatilité de l'indice
à ce strike et les volatilités des composantes au même moneyness, elle varie avec le strike~: les skews d'indice sont plus pentus que
les skews sur actions individuelles, si bien que la \omterm{def:dv:multi-asset-options:implied}{corrélation implicite} est plus élevée aux bas strikes.

\begin{definition}[Skew de corrélation]\label{def:dv:multi-asset-options:skew}
Le \emph{skew de corrélation}\index{skew de corrélation} est la variation de la \omterm{def:dv:multi-asset-options:implied}{corrélation implicite} avec le strike~: plus élevée pour
les options sur indice dont le strike est sous la monnaie, où le marché intègre des actions qui chutent ensemble, que pour celles situées au-dessus. C'est
l'équivalent actions du skew de corrélation de base des tranches de crédit (livre~2, chapitre~24).
\end{definition}

\begin{example}[Un indice synthétique de 20 actions]\label{ex:dv:multi-asset-options:index}
Prenons un indice de vingt actions à pondérations égales, aux smiles plats, dont les volatilités sont réparties régulièrement de 22~\% à 36~\%, et que
son smile à un an soit celui du chapitre~9~: 21.9~\% à 90, 19.2~\% à 100, 17.0~\% à 110. La \omterm{def:dv:multi-asset-options:implied}{corrélation implicite} vaut 0.549 à 90,
0.408 à 100 et 0.308 à 110 (la \cref{fig:dv:multi-asset-options:skew}, à gauche). Avec une corrélation constante de 0.408,
les composantes produisent un smile d'indice plat, à 19.3~\% pour les trois strikes~: une corrélation constante ne peut pas produire un
skew d'indice à partir de smiles de composantes plats. En termes de \omterm{def:dv:variance-swaps-and-volatility-derivatives:vs}{swap de variance}, avec la volatilité du strip de l'indice de 22.1~\%, la \omterm{def:dv:multi-asset-options:implied}{corrélation implicite}
vaut 0.559.
\end{example}

La \omterm{def:dv:multi-asset-options:implied}{corrélation implicite} a été, en moyenne, plus élevée que celle que les actions ont alors réalisée. Driessen, Maenhout et Vilkov
rapportent des moyennes de 39.5~\% en implicite contre 32.5~\% en réalisé pour le S\&P~500, et 46.0~\% contre 35.5~\% pour le Dow Jones
30. Ils lisent l'écart comme une large prime négative pour le risque de corrélation, et y rattachent la \omterm{def:dv:variance-swaps-and-volatility-derivatives:vrp}{prime de risque de variance} de l'indice
du chapitre~14. Le trade qui la capte porte un nom.

\begin{definition}[Trade de dispersion]\label{def:dv:multi-asset-options:dispersion}
Un \emph{trade de dispersion}\index{trade de dispersion} vend la volatilité de l'indice et achète celle des composantes, par \omterm{def:dv:variance-swaps-and-volatility-derivatives:vs}{swaps de variance}
ou straddles, en montants qui annulent l'exposition à la volatilité de chaque composante. Il reste une position courte sur
la \omterm{def:dv:multi-asset-options:implied}{corrélation implicite} de l'indice~: il gagne lorsque la corrélation réalisée passe sous la \omterm{def:dv:multi-asset-options:implied}{corrélation implicite}.
\end{definition}

\begin{definition}[Swap de corrélation]\label{def:dv:multi-asset-options:corrswap}
Un \emph{swap de corrélation}\index{swap de corrélation} verse la moyenne pondérée des corrélations deux à deux réalisées
entre les composantes d'un panier, $\sum_{i<j}w_iw_j\hat\rho_{ij}/\sum_{i<j}w_iw_j$, contre un strike fixe. Il isole la
corrélation qu'un \omterm{def:dv:multi-asset-options:dispersion}{trade de dispersion} ne détient que de manière approximative.
\end{definition}

Les pondérations du \omterm{def:dv:multi-asset-options:dispersion}{trade de dispersion} découlent de la formule de variance de l'indice. À corrélation fixée, la dérivée de la
variance de l'indice par rapport à la variance de la composante $i$ vaut
$w_i^2(1-\bar\rho)+\bar\rho\,w_i\sum_jw_j\sigma_j/\sigma_i$. Un \omterm{def:dv:variance-swaps-and-volatility-derivatives:vs}{swap de variance} court sur l'indice, de \omterm{def:dv:variance-swaps-and-volatility-derivatives:varnot}{notionnel de variance} $N$,
couvert par des \omterm{def:dv:variance-swaps-and-volatility-derivatives:vs}{swaps de variance} sur les composantes, d'un montant de $N$ fois ces dérivées, est insensible aux petites variations de la
volatilité de chaque composante. Son P\&L vaut alors $N(\bar\rho-\hat\rho)C$, où $C=(\sum w_i\sigma_i)^2-\sum w_i^2\sigma_i^2$ est le
terme croisé de la variance de l'indice.

\begin{example}[Corrélation inférieure de quinze points à l'implicite]\label{ex:dv:multi-asset-options:dispersion}
Vendre une variance d'indice à un an avec un \omterm{def:dv:variance-swaps-and-volatility-derivatives:veganot}{notionnel vega} de 100\,000, au niveau du strip de 22.12~\%, un \omterm{def:dv:variance-swaps-and-volatility-derivatives:varnot}{notionnel de variance} de 2\,261 par
point de variance, et acheter la variance des composantes selon les montants de couverture. Si chaque composante réalise sa volatilité implicite et
la corrélation réalise 0.409, quinze points sous l'implicite de 0.559, le P\&L vaut
$2\,261\times0.15\times798=270\,600$, avec $C=798$ points de variance. Sur 4\,000 années simulées de rendements quotidiens, le
P\&L moyen s'élève à 271\,000, avec un écart type de 34\,000 dû au bruit des volatilités et de la corrélation réalisées. Neuf
fois sur dix, il se situe entre 215\,000 et 326\,000 (la \cref{fig:dv:multi-asset-options:dispersion}).
\end{example}

\begin{omfigure}
\begin{tikzpicture}
\begin{groupplot}[group style={group size=2 by 1,horizontal sep=1.6cm},
  width=6.6cm,height=5cm,
  tick label style={font=\small},label style={font=\small},grid=major,grid style={gray!25},table/col sep=comma]
\nextgroupplot[xlabel={strike de l'indice},ylabel={corrélation implicite},xmin=87,xmax=113,ymin=0.25,ymax=0.6,
  xtick={90,100,110},title={\small skew de corrélation}]
\addplot[omThm,very thick,mark=*,mark size=2pt] table[x=k,y=implied_corr]{figdata/derivatives/17-multi-asset-options/corr_skew.csv};
\nextgroupplot[xlabel={log-moneyness de l'indice $x$},ylabel={corrélation locale $\rho(x)$},xmin=-0.25,xmax=0.25,ymin=0,ymax=1.05,
  title={\small corrélation locale calibrée}]
\addplot[omDef,very thick] table[x=x,y=rho]{figdata/derivatives/17-multi-asset-options/local_corr.csv};
\end{groupplot}
\end{tikzpicture}

\omcaption{Un indice synthétique de vingt actions aux smiles plats, dont le smile propre est celui du chapitre~9. À gauche~: la
\omterm{def:dv:multi-asset-options:implied}{corrélation implicite} aux trois strikes. À droite~: la \omterm{def:dv:multi-asset-options:local}{corrélation locale}, fonction du niveau de l'indice, qui reproduit le smile de l'indice
aux trois strikes. Elle est plafonnée à un sous la monnaie et diminue lorsque l'indice monte d'environ 10~\% au-dessus de son
niveau initial~; la remontée au-delà est un artefact d'une quadratique ajustée à trois strikes, qui n'imposent rien à cet endroit. Données~: le
tutoriel.}\label{fig:dv:multi-asset-options:skew}
\end{omfigure}

\begin{omfigure}
\begin{tikzpicture}
\begin{axis}[width=10.5cm,height=5cm,
  xlabel={P\&L du trade de dispersion (milliers)},ylabel={part des simulations},
  tick label style={font=\small},label style={font=\small},
  xmin=100,xmax=450,ymin=0,ymax=0.15,grid=major,grid style={gray!25},
  yticklabel style={/pgf/number format/fixed,/pgf/number format/precision=2},
  table/col sep=comma]
\addplot[omThm,very thick,const plot,fill=omThm!25] table[x=pnl,y=share]{figdata/derivatives/17-multi-asset-options/dispersion.csv} \closedcycle;
\draw[black,dashed,thick] (axis cs:270.6,0) -- (axis cs:270.6,0.15);
\node[font=\footnotesize,anchor=west] at (axis cs:335,0.12) {$N(\bar\rho-\hat\rho)C=270.6$};
\end{axis}
\end{tikzpicture}

\omcaption{Le P\&L d'un \omterm{def:dv:multi-asset-options:dispersion}{trade de dispersion} à un an (variance d'indice vendue, \omterm{def:dv:variance-swaps-and-volatility-derivatives:veganot}{notionnel vega} de 100\,000, variance des composantes achetée
aux montants de couverture) sur 4\,000 années simulées au cours desquelles la corrélation se réalise quinze points sous l'implicite. La ligne
pointillée est la formule. Données~: le tutoriel.}\label{fig:dv:multi-asset-options:dispersion}
\end{omfigure}

La dispersion n'est pas de l'argent gratuit. Le trade est court en corrélation, et le \omterm{def:dv:multi-asset-options:skew}{skew de corrélation} exprime l'idée du marché selon laquelle
la corrélation monte lorsque l'indice baisse. Dans une phase de vente, la jambe courte sur l'indice perd alors davantage que ne gagnent les jambes sur les composantes. La lecture en termes de limites à l'arbitrage de
Driessen, Maenhout et Vilkov consiste à dire que des frictions de marché réalistes rognent la prime~: le trade implique des dizaines de positions sur
options, rééquilibrées à mesure que les pondérations et les volatilités se déplacent.

\section{Corrélation locale}

L'\cref{ex:dv:multi-asset-options:index} a montré que des composantes aux smiles plats et à corrélation constante ne peuvent pas produire le
skew de l'indice. Le skew de l'indice est plus pentu que ceux des composantes, car la corrélation monte lorsque l'indice baisse. Langnau
soutient qu'une copule gaussienne (livre~4, chapitre~15) appliquée aux distributions des actions individuelles ne parvient pas à expliquer la
pente du skew de l'indice. Il a étendu la \omterm{def:dv:local-volatility:model}{volatilité locale} multi-actifs en faisant de la corrélation une fonction de l'état de
du marché.

\begin{definition}[Corrélation locale]\label{def:dv:multi-asset-options:local}
Un modèle de \emph{corrélation locale}\index{corrélation locale} attribue à chaque actif sa propre \omterm{def:dv:local-volatility:model}{volatilité locale}, calibrée sur ses propres
options, et fait de la corrélation entre eux une fonction du temps et de l'état du marché, typiquement le niveau de l'indice~:
$\rho_{ij}(t,I_t)$. La fonction est choisie de sorte que les options sur indice du modèle reproduisent le smile de l'indice, ce qu'elle fait par
construction, de la même manière que la \omterm{def:dv:local-volatility:model}{volatilité locale} ajuste le smile d'un seul actif.
\end{definition}

La calibration exacte utilise la \omterm{def:dv:local-volatility:projection}{projection markovienne} du chapitre~9. La variance locale de l'indice doit égaler l'espérance conditionnelle
de la variance instantanée du panier sachant le niveau de l'indice, et cette équation est résolue pour la
corrélation à chaque niveau de l'indice. La construction utilise la version la plus simple qui montre le mécanisme. La corrélation est
commune à toutes les paires et quadratique en le \omterm{def:dv:implied-volatility-and-its-surface:surface}{log-moneyness} de l'indice $x$, $\rho(x)=a-bx+cx^2$ bornée à $[0,1]$, et ses trois
coefficients sont calibrés sur les volatilités de l'indice aux trois strikes par des pas de Newton sur une simulation à nombres aléatoires
communs.

\begin{example}[Ajustement du skew de l'indice]\label{ex:dv:multi-asset-options:lc}
Pour l'indice synthétique, l'ajustement donne $\rho(x)=0.384-3.62x+16.3x^2$, qui reproduit l'indice à 21.93~\%, 19.19~\% et
16.99~\% contre 21.93~\%, 19.19~\% et 16.98~\% (la \cref{fig:dv:multi-asset-options:skew}, à droite). La corrélation vaut 0.91 lorsque
l'indice baisse de 10~\% en termes logarithmiques, plafonnée à un sous ce seuil, et vaut 0.19 lorsqu'il monte de 10~\%. Les lois marginales des composantes sont
inchangées, lognormales à leurs propres volatilités. Seule leur dépendance évolue avec le marché.
\end{example}

La \omterm{def:dv:multi-asset-options:local}{corrélation locale} price ensuite chaque produit sur les composantes de manière cohérente avec les deux marchés d'options. Cela inclut les \omterm{def:dv:multi-asset-options:worstof}{options worst-of},
les \omterm{def:dv:multi-asset-options:basket}{options sur panier} de strike éloigné de la monnaie et les structures de dispersion. Comme la \omterm{def:dv:local-volatility:model}{volatilité locale}, elle fixe la dynamique de
la corrélation à une fonction de l'état, et elle prévoit donc un \omterm{def:dv:multi-asset-options:skew}{skew de corrélation} forward qui s'aplatit, le pendant multi-actifs
du \omterm{def:dv:local-volatility:forward}{smile forward} du chapitre 16. Les produits worst-of dont le risque se situe dans des scénarios où l'indice a déjà chuté héritent de cette hypothèse.
baissé héritent de cette hypothèse.

\section{Options quanto et composites}

Un paiement peut faire intervenir une seconde devise même lorsqu'il se réfère à un seul actif. Les deux cas standard diffèrent selon que le
taux de change est fixé à l'avance ou converti à l'échéance.

\begin{definition}[Option quanto]\label{def:dv:multi-asset-options:quanto}
Une \emph{option quanto}\index{option quanto} verse un paiement calculé sur un actif étranger, exprimé en unités de la devise étrangère, mais
réglé en devise domestique à un taux de change fixe~: par exemple $(S_T-K)^+$ dollars par point d'indice d'un
indice en yens.
\end{definition}

\begin{definition}[Ajustement quanto]\label{def:dv:multi-asset-options:quantoadj}
L'\emph{ajustement quanto}\index{ajustement quanto} est la variation de la dérive de l'actif lorsqu'il est évalué sous la
mesure domestique avec une conversion fixe~: $-\rho\sigma_S\sigma_X$, avec $\sigma_X$ la volatilité du taux de change
(unités domestiques par unité étrangère) et $\rho$ sa corrélation avec l'actif. Le forward quanto vaut $Fe^{-\rho\sigma_S\sigma_XT}$.
\end{definition}

L'ajustement découle du théorème de Girsanov (livre~4, chapitre~5). Sous la mesure étrangère, la dérive de l'actif
est le carry étranger. Un investisseur domestique qui détient l'actif converti au taux flottant perçoit le taux domestique sur
$S_tX_t$, ce qui fixe la dérive de $SX$. En retranchant la dérive propre du taux de change et le terme de covariance
$\rho\sigma_S\sigma_X$ reste la dérive de l'actif sous la mesure domestique. Le hedger d'une \omterm{def:dv:multi-asset-options:quanto}{option quanto} détient l'actif étranger
et doit continuellement réajuster la couverture de sa valeur en devise étrangère. Lorsque l'actif et la devise sont corrélés, ces
couvertures gagnent ou perdent systématiquement, et l'ajustement en tient compte.

\begin{example}[Un forward sur indice payé dans une autre devise]\label{ex:dv:multi-asset-options:quanto}
Un indice étranger dont le forward vaut 100, de volatilité 20~\%, payé à un an à raison d'une unité domestique par point. Avec une
volatilité du taux de change de 10~\% et des corrélations de $-0.6$, $-0.3$, $0.3$ et $0.6$, le forward quanto vaut 101.21, 100.60,
99.40 et 98.81. Une simulation de 200\,000 trajectoires sous la mesure domestique avec la dérive ajustée concorde à 0.005
(la \cref{fig:dv:multi-asset-options:quanto}). Pour une corrélation de $-0.3$, le call quanto à la monnaie vaut 8.29 contre
7.97 pour le même call réglé en devise étrangère.
\end{example}

\begin{definition}[Option composite]\label{def:dv:multi-asset-options:composite}
Une \emph{option composite}\index{option composite} est une option sur la valeur de l'actif étranger convertie dans la
devise domestique au taux en vigueur, $(S_TX_T-K)^+$ avec un prix d'exercice domestique. Son sous-jacent $SX$ a une volatilité
$\sqrt{\sigma_S^2+\sigma_X^2+2\rho\sigma_S\sigma_X}$.
\end{definition}

Une \omterm{def:dv:multi-asset-options:quanto}{option quanto} retire le risque de change du paiement et laisse la corrélation dans la dérive. Une \omterm{def:dv:multi-asset-options:composite}{option composite} conserve le
risque de change et place la corrélation dans la volatilité~: pour $\rho=-0.3$, la volatilité composite vaut $\sqrt{0.04+0.01-0.012}=19.5\%$.
Ce sont deux options à deux actifs, et toutes deux portent la même position de corrélation, que le desk doit évaluer et couvrir. Les options spread
entre deux actifs, l'autre produit courant à deux actifs, sont traitées dans le livre~6, chapitre~16.

\begin{omfigure}
\begin{tikzpicture}
\begin{axis}[width=10.5cm,height=5cm,
  xlabel={corrélation entre l'indice et le taux de change},ylabel={forward quanto},
  tick label style={font=\small},label style={font=\small},
  xmin=-0.7,xmax=0.7,ymin=98.5,ymax=101.5,grid=major,grid style={gray!25},
  legend columns=2,legend style={font=\footnotesize,at={(0.5,-0.3)},anchor=north,draw=none,fill=none,/tikz/every even column/.append style={column sep=8pt}},
  table/col sep=comma]
\addplot[omThm,very thick] table[x=rho,y=analytic]{figdata/derivatives/17-multi-asset-options/quanto.csv};
\addplot[black,only marks,mark=*,mark size=1.8pt] table[x=rho,y=mc]{figdata/derivatives/17-multi-asset-options/quanto.csv};
\legend{$Fe^{-\rho\sigma_S\sigma_XT}$,simulation}
\end{axis}
\end{tikzpicture}

\omcaption{Le forward à un an d'un indice étranger (forward 100, volatilité 20~\%) payé à un taux de change fixe dont la
volatilité est de 10~\%, en fonction de la corrélation entre l'indice et le taux. Données~: le tutoriel.}\label{fig:dv:multi-asset-options:quanto}
\end{omfigure}

\section{Tutoriel~: corrélation en pratique}\label{tut:dv:multi-asset-options:tut}

\begin{tutorial}
\textbf{Objectif.} Évaluer les paniers, les \omterm{def:dv:multi-asset-options:worstof}{options worst-of} et best-of en fonction de la corrélation~; lire la \omterm{def:dv:multi-asset-options:implied}{corrélation implicite} d'un indice à trois
strikes~; calibrer une \omterm{def:dv:multi-asset-options:local}{corrélation locale} sur le skew d'indice~; simuler un \omterm{def:dv:multi-asset-options:dispersion}{trade de dispersion}~; mesurer l'\omterm{def:dv:multi-asset-options:quantoadj}{ajustement quanto}.
\textbf{État final~:} les quatre figures et les chiffres du problème du week-end.
\begin{tutsteps}
\item \textbf{Trajectoires en \omterm{def:dv:multi-asset-options:local}{corrélation locale}.} Chocs équicorrélés dont la corrélation commune se lit sur le niveau de l'indice à chaque
pas~:
\omcode{code/firm/multiasset/firm_multiasset.py}{51}{75}{Le générateur de trajectoires en corrélation locale.}\label{lst:dv:multi-asset-options:lc}
\item \textbf{\omterm{def:dv:multi-asset-options:implied}{Corrélation implicite} et variance de l'indice}, les deux sens d'une même formule~:
\omcode{code/firm/multiasset/firm_multiasset.py}{119}{130}{Corrélation implicite à partir des volatilités de l'indice et des composantes.}\label{lst:dv:multi-asset-options:implied}
\item \textbf{Le forward quanto et la volatilité composite}~:
\omcode{code/firm/multiasset/firm_multiasset.py}{143}{151}{Ajustements quanto et composite.}\label{lst:dv:multi-asset-options:quanto}
\item \textbf{Exécuter} \texttt{dv\_multiasset.three\_share\_prices()}, \texttt{correlation\_by\_strike()},
\texttt{calibrate\_local\_correlation()}, \texttt{dispersion()}, \texttt{quanto\_example()} et \texttt{fig\_multiasset.py}.
\end{tutsteps}
\textbf{Pour aller plus loin.} Donner à chaque composante son propre skew (une \omterm{def:dv:local-volatility:model}{volatilité locale} chacune) et recalibrer la \omterm{def:dv:multi-asset-options:local}{corrélation locale}~;
évaluer le digital worst-of à trois actions sous la \omterm{def:dv:multi-asset-options:local}{corrélation locale} calibrée~; exécuter le \omterm{def:dv:multi-asset-options:dispersion}{trade de dispersion}, la corrélation
se réalisant au niveau implicite, mais avec des volatilités de composantes inférieures de cinq points.
\end{tutorial}

\section{Construction~: trajectoires et paiements multi-actifs}\label{bld:dv:multi-asset-options:multiasset}

\begin{build}
\bfield{Rôle} Le moteur multi-actifs de la firme miniature~: trajectoires corrélées pour les paniers, les \omterm{def:dv:multi-asset-options:worstof}{options worst-of} et quanto, les outils d'analyse de la corrélation implicite-
et réalisée de son book de dispersion, et le générateur sur lequel reposent les autocallables du chapitre 18.
\bfield{Interface} \texttt{correlated\_paths(s0, vols, corr, times, r, q, n\_paths, seed, quanto, steps\_per\_date)} avec
\texttt{vols[i]} un nombre ou une fonction $\sigma_i(t,S)$~; \texttt{local\_correlation\_paths(s0, vols, weights, rho\_fn, times,
\ldots)}~; \texttt{equicorrelation}~; \texttt{basket\_payoff}, \texttt{worst\_of\_payoff}, \texttt{best\_of\_payoff},
\texttt{worst\_of\_digital}~; \texttt{basket\_moment\_match}~; \texttt{implied\_correlation}, \texttt{index\_variance},
\texttt{realised\_correlation}~; \texttt{quanto\_forward}, \texttt{composite\_vol}.
\bfield{Règles} Les matrices de corrélation sont vérifiées définies positives avant usage (la factorisation de Cholesky échoue sinon)~; la
corrélation quanto est cotée par rapport au taux en unités domestiques par unité étrangère~; chaque nombre de Monte-Carlo porte une erreur type.
\bfield{Tests d'acceptation} \texttt{code/firm/multiasset/tests/}~: les actifs simulés sont des martingales avec la corrélation spécifiée~;
les \omterm{def:dv:multi-asset-options:worstof}{options worst-of} et best-of se réduisent à l'option vanille à corrélation un~; l'appariement des moments par rapport à la simulation~; la \omterm{def:dv:multi-asset-options:implied}{corrélation implicite} inverse
la variance de l'indice~; la \omterm{def:dv:multi-asset-options:local}{corrélation locale} se réduit au cas constant~; le forward quanto concorde avec une simulation.
\bfield{Pour aller plus loin} \omterm{def:dv:local-volatility:model}{Volatilités locales} par composante et \omterm{def:dv:multi-asset-options:local}{corrélation locale} par paires calibrée par \omterm{def:dv:local-volatility:projection}{projection markovienne}~; le
rapport cross-gamma du livre~6.
\end{build}

\begin{omsources}
\item J.~Driessen, P.~Maenhout et G.~Vilkov, ``The price of correlation risk: evidence from equity options'', SSRN 673425
(2005)~; ``Option-implied correlations and the price of correlation risk'', SSRN 2359380 (2013).
\item A.~Langnau, ``Introduction into `local correlation modelling''', arXiv 0909.3441 (2009).
\end{omsources}

\section{Exercices}

\begin{exercise}[$\star$]\label{exo:dv:multi-asset-options:1}
Deux actifs de volatilités 20~\% et 30~\%, pondérés à parts égales, de corrélation 0.5~: quelle est la volatilité du panier dans le cadre de l'approximation
lognormale~?
\end{exercise}

\begin{exercise}[$\star$]\label{exo:dv:multi-asset-options:2}
Pourquoi une option d'achat best-of perd-elle de la valeur lorsque la corrélation augmente, alors qu'une option d'achat sur panier en gagne~?
\end{exercise}

\begin{exercise}[$\star$]\label{exo:dv:multi-asset-options:3}
Un indice de deux composantes à parts égales, de volatilités 25~\% et 35~\%, se traite à 24~\%. Calculer sa \omterm{def:dv:multi-asset-options:implied}{corrélation implicite}.
\end{exercise}

\begin{exercise}[$\star\star$]\label{exo:dv:multi-asset-options:4}
Établir le forward quanto $Fe^{-\rho\sigma_S\sigma_XT}$ à partir de l'exigence que $S_TX_T$, converti au taux flottant,
rapporte le taux domestique.
\end{exercise}

\begin{exercise}[$\star\star$]\label{exo:dv:multi-asset-options:5}
Expliquer pourquoi un dealer qui vend des notes worst-of est vendeur de corrélation, et citer un instrument coté qui compense une partie de la
position.
\end{exercise}

\begin{exercise}[$\star\star$]\label{exo:dv:multi-asset-options:6}
À l'aide de l'indice synthétique, calculer le P\&L de dispersion si la corrélation réalisée se fixe au niveau implicite, mais que la
volatilité de chaque composante se réalise un point au-dessus de son niveau implicite (montants de couverture inchangés).
\end{exercise}

\begin{exercise}[$\star\star\star$]\label{exo:dv:multi-asset-options:7}
\emph{Programmation.} Pricer le digital worst-of sur trois actions (toutes au-dessus de 70~\%) sous la fonction de \omterm{def:dv:multi-asset-options:local}{corrélation locale} calibrée de l'indice synthétique
appliquée au panier propre aux trois actions. Comparer avec les prix à corrélation constante du chapitre.
\end{exercise}

\begin{exercise}[$\star\star\star$]\label{exo:dv:multi-asset-options:8}
\emph{Repérer l'erreur.} «~La \omterm{def:dv:multi-asset-options:implied}{corrélation implicite} a dépassé en moyenne la corrélation réalisée, donc un desk de dispersion qui vend la variance de l'indice
et achète la variance des composantes est un trade de carry sans risque.~»
\end{exercise}

\section{Problème~: l'indice et ses composantes}

\begin{problem}[{Problème du week-end --- corrélation implicite et trade de dispersion}]\label{pb:dv:multi-asset-options:1}
Un indice synthétique regroupe vingt actions à poids égaux, avec des smiles plats et des volatilités régulièrement réparties de 22~\% à 36~\%. Le
smile à un an de l'indice lui-même est celui du chapitre~9. Un desk de dispersion envisage de vendre la variance de l'indice contre la variance des composantes.

\medskip\noindent\textbf{Partie I --- Les composantes.}
\begin{enumerate}
\item Calculer $\sum w_i^2\sigma_i^2$ et le terme croisé $C$.
\item Quelle volatilité de l'indice impliquerait une corrélation nulle~? Une corrélation égale à un~?
\item Pourquoi la volatilité de l'indice est-elle inférieure à la volatilité moyenne des composantes~?
\item Donner la volatilité de l'indice à 90, 100 et 110.
\item Donner la volatilité de l'indice que produit une corrélation constante de 0.408 à 90 et à 110.
\end{enumerate}

\medskip\noindent\textbf{Partie II --- \omterm{def:dv:multi-asset-options:implied}{Corrélation implicite}.}
\begin{enumerate}[resume]
\item Donner la \omterm{def:dv:multi-asset-options:implied}{corrélation implicite} à 90, 100 et 110.
\item Donner la \omterm{def:dv:multi-asset-options:implied}{corrélation implicite} du \omterm{def:dv:variance-swaps-and-volatility-derivatives:vs}{swap de variance}.
\item Expliquer le \omterm{def:dv:multi-asset-options:skew}{skew de corrélation} en une phrase.
\item Donner la fonction de \omterm{def:dv:multi-asset-options:local}{corrélation locale} qui reprice les trois strikes.
\item Quelle \omterm{def:dv:multi-asset-options:local}{corrélation locale} celle-ci attribue-t-elle à 10~\% à la baisse et à 10~\% à la hausse~?
\end{enumerate}

\medskip\noindent\textbf{Partie III --- Le \omterm{def:dv:multi-asset-options:dispersion}{trade de dispersion}.}
\begin{enumerate}[resume]
\item Pour un notionnel de \omterm{def:dv:greeks-and-the-hedging-pnl:greeks}{vega} de 100\,000 sur l'indice, donner le \omterm{def:dv:variance-swaps-and-volatility-derivatives:varnot}{notionnel de variance}.
\item Donner la formule des montants de couverture des composantes et expliquer ce qu'ils compensent.
\item Donner le P\&L si la corrélation se réalise quinze points en dessous du niveau implicite.
\item Donner la moyenne simulée, l'écart type et l'intervalle à 5--95~\%.
\item Que paierait un \omterm{def:dv:multi-asset-options:corrswap}{swap de corrélation} fixé au niveau implicite de 0.559, par unité de notionnel~?
\end{enumerate}

\medskip\noindent\textbf{Partie IV --- Jugement.}
\begin{enumerate}[resume]
\item Quand le \omterm{def:dv:multi-asset-options:dispersion}{trade de dispersion} perd-il, et de combien peut-il perdre~?
\item Que disent les travaux publiés sur la \omterm{def:dv:multi-asset-options:implied}{corrélation implicite} comparée à la corrélation réalisée~?
\item Pourquoi la prime peut-elle persister~?
\item Énoncer le \emph{résultat nommé}~: la \omterm{def:dv:multi-asset-options:implied}{corrélation implicite} de l'indice synthétique à trois strikes, et le P\&L du \omterm{def:dv:multi-asset-options:dispersion}{trade de dispersion}
lorsque la corrélation réalisée se fixe quinze points en dessous du niveau implicite.
\item En une phrase~: que price une option sur indice que les options sur ses composantes ne pricent pas~?
\end{enumerate}
\end{problem}

\section{Questions d'entretien}

\begin{interviewq}[$\star$ \iqroles{trader, chercheur}]\label{iq:dv:multi-asset-options:1}
Comment la corrélation affecte-t-elle les prix d'un call sur panier, d'un put worst-of et d'un call best-of~?
\end{interviewq}

\begin{interviewq}[$\star$ \iqroles{chercheur}]\label{iq:dv:multi-asset-options:2}
Définissez la \omterm{def:dv:multi-asset-options:implied}{corrélation implicite}. Pourquoi est-elle plus élevée pour les bas strikes~?
\end{interviewq}

\begin{interviewq}[$\star\star$ \iqroles{trader}]\label{iq:dv:multi-asset-options:3}
Expliquez un \omterm{def:dv:multi-asset-options:dispersion}{trade de dispersion}~: les jambes, les pondérations et ce à quoi vous êtes exposé.
\end{interviewq}

\begin{interviewq}[$\star\star$ \iqroles{chercheur, développeur}]\label{iq:dv:multi-asset-options:4}
Établissez l'\omterm{def:dv:multi-asset-options:quantoadj}{ajustement quanto}. Quel est son signe pour un indice japonais payé en dollars si l'indice tend à baisser quand le yen
se renforce~?
\end{interviewq}

\begin{interviewq}[$\star\star$ \iqroles{développeur}]\label{iq:dv:multi-asset-options:5}
Votre matrice de corrélation pour 30 actifs échoue à la factorisation de Cholesky. Pourquoi, et que faites-vous~?
\end{interviewq}

\begin{interviewq}[$\star\star\star$ \iqroles{risque, trader}]\label{iq:dv:multi-asset-options:6}
Un livre d'autocallables worst-of est vendeur de corrélation. Comment mesureriez-vous, couvririez-vous et provisionneriez-vous cette exposition~?
\end{interviewq}
